%% Livret de séquence — généré pour rendu agrégé.
%% Niveau : N11 — Séquence : S01 — Représentations d'un nombre

\documentclass[a4paper,11pt]{article}

%% ── Paquets externes (noyau pour atome) ───────────────
\usepackage[utf8]{inputenc}
\usepackage[T1]{fontenc}
\usepackage{lmodern}
\usepackage[french]{babel}
\usepackage{geometry}
\usepackage{amsmath}
\usepackage{amssymb}
\usepackage{etoolbox}
\usepackage{xkeyval}
\usepackage{xifthen}
\usepackage{environ}
\usepackage{forloop}
\usepackage{xcolor}
\usepackage{ninecolors}
\usepackage{multirow}
\usepackage{colortbl}
\usepackage{tabularray}
\usepackage{multicol}
\usepackage[inline]{enumitem}
\usepackage{graphicx}
\usepackage{adjustbox}
\usepackage{tikz}
\usepackage[theorems,breakable,skins]{tcolorbox}
\usepackage{pifont}

%% ── Marges (alignées sur les livrets) ───────────────
\newgeometry{left=35pt,right=35pt,top=65pt,bottom=70pt}

%% ── Paquets externes additionnels (selon l'atome) ─────
\usepackage{datetime}
\usepackage{tabularx}

%% ── Initialisations de bibliothèques ─────────────────
\UseTblrLibrary{booktabs}
\UseTblrLibrary{varwidth}

%% ── Définitions du paquet seqenseigne (inlinées) ─────
%%    Fermeture transitive des macros effectivement utilisées
%%    par l'atome — au lieu de \usepackage{seqenseigne}
%%    qui chargerait ~110 paquets.
\makeatletter
\makeatother

%% ── Macros transverses ───────────────────────────────
\providecommand{\acompleter}[1]{\;\underline{\phantom{#1}}\;}

\usepackage{titlesec}
\usepackage{fancyhdr}
\usepackage{lastpage}

\geometry{vmargin=70pt,hmargin=50pt,headheight=50pt,headsep=15pt,footskip=20pt}

\pagestyle{fancy}
\fancyhf{}
\renewcommand{\headrulewidth}{0pt}
\rhead{\rightmark}
\rfoot{Page \thepage/\pageref{LastPage}}
\lfoot{N11 — S01}

\graphicspath{{data/images/}{./images/}{./}}

\newdateformat{datemoisannee}{\monthname\;\THEYEAR}

\begin{document}

\providecommand{\theHExoNum}{\theExoNum}
\providecommand{\theHNotionNum}{\theNotionNum}
\providecommand{\theHMethodeNum}{\theMethodeNum}

\seqCreeCompteurs
\seqSetCodeNiveau{N11}
\seqSetCodeSequence{S01}
\seqSetColorsTheme{nombres}
\seqInitCorriges
\seqInitAnnexes
\seqInitRemediation

%% Boîte de titre (faite à la main).
\begin{boiteTitreGen}{Séquence 01}
  \begin{minipage}[c][100pt][c]{\linewidth}
    \vfill \textit{Livret de séquence} \hfill
    \vfill \centering {\Huge \textbf{Représentations d'un nombre}}
    \vfill Version de \datemoisannee\today \hfill \textit{Classe de 4\ieme}
  \end{minipage}
\end{boiteTitreGen}

%% Objectifs générés depuis la BDD — v0.51.17 : une boîte par partie, en page 1.
\seqRenduTableauObjectifs{%
  \textbf{01.} Connaître le cours &  &  &  \\\hline
  \textbf{02.} Calculer une fréquence &  &  & 
}

\clearpage
\begin{seqSerieExos}{0}
\begin{seqExercice}[nom={EA1}]
Énoncé EA1
\seqCorrige{%
Corrigé EA1
}
\end{seqExercice}
\end{seqSerieExos}

\clearpage
\seqTitreSection{Connaissances}

\begin{seqNotion}{Notion 1}{La fréquence est ...}
\end{seqNotion}

\clearpage
\seqTitreSection{Savoir-faire}

\begin{seqMethode}{Calculer une fréquence}{On compte ...}
\end{seqMethode}

\clearpage
\begin{seqSerieExos}{1}
\begin{seqExercice}[nom={F1},obj={02},remediation=oui]
Énoncé F1
\seqCorrige{%
Corrigé F1
}
\seqRemediation{%
Remède F1
}{%
Corr R F1
}
\end{seqExercice}
\end{seqSerieExos}

\clearpage
\begin{seqSerieExos}{2}
\begin{seqExercice}[nom={A1},obj={02}]
Énoncé A1
\seqCorrige{%
Corrigé A1
}
\end{seqExercice}
\end{seqSerieExos}

\clearpage
\begin{seqSerieExos}{3}
\begin{seqExercice}[nom={E1},obj={02}]
Énoncé E1
\seqCorrige{%
Corrigé E1
}
\end{seqExercice}
\end{seqSerieExos}

\clearpage
\seqAfficheAnnexes
\clearpage
\seqAfficheCorriges

\seqAfficheRemediations
\seqAfficheCorrigesRemediation

\end{document}